Tomographic measurements of gravitational lensing with different lens and source redshift distribution contain crucial information about the relative expansion rate of the Universe, and hence dark energy.
%
% MSD
While this technique is well-established in weak lensing, its application to the strong lensing regime has traditionally focused on rare Double Source Plane Lenses (DSPLs). However, the cosmological utility of DSPL systems is fundamentally limited by the Mass-Sheet Degeneracy (MSD)—a systematic uncertainty that has been underexplored in previous forecasting literature—as well as their intrinsic scarcity.
%
In this work, we address these challenges through two novel contributions. First, we incorporate the MSD into a hierarchical forecasting framework, demonstrating that this degeneracy severely degrades constraints from small samples of genuine DSPLs.
%
% PMSPLs
Second, to overcome the intrinsic rarity of DSPLs and leverage population statistics to break these degeneracies, we introduce the concept of \emph{Pseudo} Double-Source Plane Lenses (PDSPLs). We generalize the double-source plane formalism to pairs of independent single-source plane lenses with self-similar deflectors, unlocking the statistical power of the $\sim\mathcal{O}(10^5)$ systems expected from upcoming surveys like LSST, Euclid and Roman.
%
% Results
We derive the transformation of the Einstein radius ratio under the degrees of freedom of an internal mass sheet and apply this to our forecasted sample. We demonstrate that the PDSPL methodology can yield tight constraints on the dark energy equation of state parameter $w_0$ when inferred under a Flat $w_0w_a$CDM cosmology, achieving a precision of $\sigma(w_0) \sim 0.43$ with the LSST 10-year photometric galaxy-galaxy lens sample alone. Notably this precision is achieved while simultaneously constraining MSD and deflector power law slope to $\sim2\%$. This constraint improves to $\sigma(w_0) \sim 0.27$ when combined with a prior $\mathcal{N}(0.3, 0.05)$ on the matter density $\Omega_{\rm m}$, effectively turning strong lensing tomography into a robust statistical probe of cosmic expansion.
%
% WL comparison
% Our constraints are comparable with those from current weak lensing analysis (DES Y3 $3\times 2$pt).